\documentclass[10pt,a4paper]{article}
\usepackage[T1]{fontenc}
\usepackage[utf8]{inputenc}
\usepackage{lmodern}
\usepackage[margin=1.8cm]{geometry}
\usepackage{amsmath,amssymb,booktabs,tabularx,enumitem}
\usepackage[hidelinks]{hyperref}
\hypersetup{pdftitle={E-BBB-3: Simulation study of the framework}}
\usepackage{microtype,fancyhdr}
\setlength{\parindent}{0pt}
\setlength{\parskip}{4pt}
\setlist{nosep,leftmargin=1.35em}
\pagestyle{fancy}
\fancyhf{}
\fancyhead[L]{\small E-BBB-3}
\fancyhead[R]{\small Simulation study of the framework}
\fancyfoot[C]{\small\thepage}
\setlength{\headheight}{13pt}
\newcommand{\E}{\mathbb E}
\newcommand{\Prob}{\mathbb P}
\newcommand{\logit}{\operatorname{logit}}
\newcommand{\HN}{\operatorname{HN}}
\begin{document}
\begin{center}
{\Large\bfseries E-BBB-3: simulation study of the framework}\\[4pt]
{\small Research and software set-up \enspace|\enspace 26 November 2026}
\end{center}
\vspace{-0.5em}

\section*{1. Aim and scope}
Evaluate whether Bayesian borrowing improves the power and efficiency of
anytime-valid tests in a basket trial with binary responses. The lead
document is \texttt{framework.tex}. This first study compares borrowing
choices within \emph{one} e-process. Design-optimal e-values, Bayesian
decision rules, Simon designs, flexible inclusion and new entry-time
redesigns remain outside this study in \texttt{Additional\_files/legacy/}.

A basket is a cell $c=(j,l)$ of a tumour-type $\times$ biomarker grid.
These labels are a working application, to be confirmed for a clinical
protocol. The current simulation has independent Bernoulli outcomes with
constant response probability $p_c$, deterministic recruitment and immediate
outcome observation. It implements a useful special case of the broader
conditional framework; no clinical data or external data $D_0$ are used.

\section*{2. Test and guarantee}
Let $\mathcal G_i$ contain all observed platform data before outcome $Y_i$
and the current basket label $c_i$. For each basket separately, test
\[
 H_{0,c}:\quad \Prob(Y_i=1\mid\mathcal G_i)\le p_0
 \quad\hbox{for every }i\hbox{ assigned to }c.
\]
This is a conditional null, not merely a restriction on the average response
rate. In the simulated independent, constant-rate setting it reduces to
$p_c\le p_0$.

At analysis $\tau_{c,k-1}$, fit the borrowing model to all platform data
available then, including the basket's own earlier outcomes. Restrict its
posterior for $p_c$ to $(p_0,1]$ and normalize:
\[
 \pi_{c,k-1}(A)=
 \frac{\widetilde\pi_{c,k-1}(A\cap(p_0,1])}
      {\widetilde\pi_{c,k-1}((p_0,1])}.
\]
Freeze this probability measure before the next block. For the new block
of $n_{c,k}$ outcomes and $y_{c,k}$ responses, use
\[
 e_{c,k}=\int
 \left(\frac p{p_0}\right)^{y_{c,k}}
 \left(\frac{1-p}{1-p_0}\right)^{n_{c,k}-y_{c,k}}
 \pi_{c,k-1}(dp),\qquad
 E_{c,k}=\prod_{r=1}^k e_{c,r},\quad E_{c,0}=1.
\]
Reject at the first scheduled analysis with $E_{c,k}\ge 1/\alpha$.

\textbf{Why it is valid.} For fixed $p>p_0$ and
$q_i=\Prob(Y_i=1\mid\mathcal G_i)\le p_0$,
\[
 \E\!\left[
  (p/p_0)^{Y_i}\{(1-p)/(1-p_0)\}^{1-Y_i}
  \mid\mathcal G_i\right]
 =q_i\,p/p_0+(1-q_i)(1-p)/(1-p_0)\le1.
\]
Products within a block, predictable probability mixtures and optional
sampling preserve the non-negative supermartingale argument. Ville's
inequality gives $\Prob_{H_{0,c}}(\sup_k E_{c,k}\ge1/\alpha)\le\alpha$.
Borrowing-model correctness is unnecessary. The supplied prior must actually
be normalized, supported above $p_0$ and frozen; numerical fitting randomness
must not reveal future outcomes.

The guarantee is \textbf{per basket}. There is no platform-wide family-wise
error rate (FWER) or false discovery rate guarantee. Outcome-dependent
delays, allocation and analysis rules require their own specification under
the conditional null; they are not evaluated in this first simulation.

\newpage
\section*{3. Borrowing models and numerical implementation}
The priors follow the eBBB code of T.~Jacobs (\texttt{LR\_evalue\_multiple}),
with our two-way EX/NEX in place of its Stan fit. Each is fitted to all
platform counts $(y_k,n_k)$ at the analysis, then restricted and frozen as in
Section~2. Beta-based priors start from $\operatorname{Beta}(a_0,b_0)$,
$a_0=b_0=1$; $R_c$ and $C_c$ are the other baskets in the same row (tumour type)
and column (biomarker group).

\begin{tabularx}{\linewidth}{@{}p{3.3cm}X@{}}
\toprule
Prior & Definition\\
\midrule
No borrowing & $\operatorname{Beta}(a_0+y_c,\ b_0+n_c-y_c)$.\\
Power, equal / column & Add $w_r\sum_{R_c}$ and $w_c\sum_{C_c}$ of the responses and non-responses; $(w_r,w_c)=(0.5,0.5)$ or $(0,1)$.\\
EX/NEX & Mixture of EX $=\operatorname{Beta}(a_0,b_0)$ plus the column data and NEX $=\operatorname{Beta}(a_0,b_0)$, each updated with the own data; $\Prob(EX)=0.5$ updated by the beta-binomial likelihood of the own data.\\
Similarity & Add every basket with data with weight $s=(1-\mathrm{JSD})^{2}$, set to 0 if $s<0.3$; JSD (base 2) between the own-data posteriors.\\
Two-way EX/NEX & Model below, hyperparameters of eBBB (\texttt{exnex.stan}).\\
Oracle & Point mass at $p_c$ if $p_c>p_0$, otherwise at $p_1$: a truth-informed reference.\\
\bottomrule
\end{tabularx}

\textbf{Two-way EX/NEX.} With $\theta_c=\logit(p_c)$ and, independently,
$\delta_c\sim\operatorname{Bernoulli}(0.5)$:
\[
 \delta_{jl}=1:\ \theta_{jl}=\mu+a_j+b_l+\varepsilon_{jl},\quad
 a_j\sim N(0,\sigma_a^2),\ b_l\sim N(0,\sigma_b^2),\ \varepsilon_{jl}\sim N(0,\sigma_e^2);
 \qquad \delta_c=0:\ \theta_c\sim N(\logit p_0,2^2),
\]
with $\mu\sim N(\logit p_0,2^2)$ on $[-10,6]$, $\sigma_a,\sigma_b\sim\HN(1)$ and
$\sigma_e\sim\HN(0.5)$; no sum-to-zero constraint. Missing cells are excluded;
included baskets without outcomes still receive borrowing.

\textbf{Computation.} The $\theta$ grid has width $0.05$ on $[-10,6]$, with
$\logit(p_0)$ at a cell boundary; the e-process sums normalized weights at the
support points $p_g>p_0$ directly. A Beta prior is put on the grid as its
density at the cell midpoints times the cell widths (tails into the end cells);
its block factors agree with the closed form of eBBB within $3\cdot10^{-4}$
(relative) in the tested states. The oracle's points are added to the grid, so it is exact. Two-way
fitting uses Metropolis-within-Gibbs with a Rao--Blackwellised grid
distribution (two chains, 500 burn-in and 1,000 retained draws), SD grids,
interpolation and a coarser reconstruction lattice; its accuracy affects
borrowing and power, not validity. The earlier one-way EX/NEX models remain in
the code as options.

\section*{4. Simulation design already defined in the code}
\begin{tabularx}{\linewidth}{@{}p{3.3cm}X@{}}
\toprule
Setting & Current implementation\\
\midrule
Hypotheses, horizon & $p_0=0.10$, reference alternative $p_1=0.30$, $\alpha=0.05$ per basket; at most 30 patients per basket.\\
Grid & $J=4$ tumour types $\times$ $L=3$ biomarker groups, as in eBBB.\\
Balanced recruitment & 12 cells, one patient per time unit per basket, all open at time zero.\\
Unbalanced recruitment & 11 cells, $(3,3)$ absent; rates 1, 0.5 or 0.25, falling with tumour type and biomarker group; basket $(2,3)$ opens at time 20.\\
Analyses & Global analyses every five time units; block sizes can differ across baskets.\\
Stopping & Reject at $E\ge20$, otherwise continue to 30 patients (efficacy or $n_{\max}$, as in eBBB).\\
\bottomrule
\end{tabularx}

Eleven response patterns: the four of eBBB (global null, biomarker~1 active,
interaction: biomarker~1 active except tumour type~4, all active) and tumour
type~1 active, one active basket, tumour type~1 and biomarker~1 active,
additive logit effects, diagonal interaction, mixed active/null cells and
strong heterogeneity. With two recruitment designs and seven priors this gives
154 study cells. Methods share potential patient outcomes within a run. At each
analysis new evidence uses the previous frozen prior; all observed counts,
including stopped baskets, then inform the next prior.

\emph{Differences from eBBB that remain by design:} calendar-time recruitment
with global analyses (eBBB: one accrual stream of 800 patients, an analysis
every five outcomes of a basket), $p_0,p_1=0.10,0.30$ (eBBB $0.30,0.50$),
$n_{\max}=30$ (eBBB 100), Beta priors on the grid rather than in closed form,
and the two-way model by our grid sampler rather than Stan draws.

\emph{Optional open-ended variant} (\texttt{simulate.R open}): no maximum per
basket; stop at $E\ge20$ or for futility at $E^f\ge1/\beta=5$, where $E^f$
tests $p\ge p_1$ with the posterior restricted to $[0,p_1)$ (framework.tex,
Proposition~2; the oracle bets on $p_c$ if $p_c<p_1$, otherwise $p_0$); 150
patients only end the computation. Details:
\texttt{Additional\_files/docs/open-ended-design.md}.

\newpage
\section*{5. Outputs and interpretation}
\texttt{simulate.R} runs the study with the settings in \texttt{setting.R};
\texttt{additional\_functions.R} contains the e-value (Part 1), the borrowing
models (Part 2), the simulation (Part 3) and the summaries and plots (Part 4).
Reference quadrature and JAGS calculations are validation tools in
\texttt{Additional\_files/validation/}, not study comparators.

The study writes per-basket rejection rates and their Monte Carlo standard
errors, expected sample size and decision time, signed posterior bias,
empirical 90\% posterior-interval coverage, posterior EX probability and
restricted-prior means. Platform summaries include FWER, average power among
active baskets, probability of any active rejection, probability all decisions
are correct, expected total sample size and duration. Mean absolute
estimation error is labelled separately from bias.
Posterior coverage and EX probabilities are diagnostics, not additional
frequentist guarantees. A dedicated \emph{false borrowing} measure is not yet
implemented.

Report each null basket's type I error separately; FWER can exceed $\alpha$.
The observed maximum over null baskets is a descriptive summary with its own
Monte Carlo variability. For 1,000 independent replications, a rejection rate
of 0.05 has standard error about $0.0069$; small differences cannot support
strong comparative claims. Outcome and fitting seeds come from disjoint
ranges per simulation cell, and every cell simulates all its trials in one
process, so results do not depend on the number of cores. Record seeds,
code version and session details.

\section*{6. Review and validation status}
\textbf{Corrections and regression coverage.}
The earlier initialization defect is corrected: prior rows are copied directly
across trials and checked for normalization. Regressions compare initial
factors with direct mixtures, check prior means and compare batched with
individual trials. The additive scenario's zero-effect basket is now set
explicitly to $p_0$: a logit round trip can round just above $p_0$ and
incorrectly label it active. Two checks preserve this exact boundary and the
single null basket in both recruitment designs.

\textbf{Independently checked evidence.}
GitHub Actions run \href{https://github.com/HenkvdPol/E-BBB-3/actions/runs/45855905668}{45855905668}
for correction commit \texttt{1f7dd90} parsed 41 R files and reported
\textbf{132 passes, zero failures and one skipped slow-reference group}.
The checks include all seven current models, fixed stopping and the earlier
batching regressions. No local R execution was available for this review;
the R execution evidence comes from the recorded CI log.

The exact-path tests enumerate $6^6$ block-count paths for one null basket
under selected other-basket histories. Rejection probabilities are
0.0117--0.0220 for the eBBB priors and 0.0059--0.0073 for the earlier engine
checks, at $\alpha=0.05$. These finite examples check implementation;
the theorem provides the general validity argument.

The earlier posterior validation covers 22 states. A new comparison covers
six states on the current $4\times3$ layout with the eBBB hyperparameters.
Independent recomputation reproduced its saved summaries: production versus
JAGS restricted-mean differences have median 0.0012 and maximum 0.0077;
relative differences in a 5/5 block factor have median 3.3\% and maximum 17\%.
These comparisons support further evaluation, not unbiased MCMC or exact
continuous-posterior computation. \href{https://github.com/HenkvdPol/E-BBB-3/pull/8}{PR\,8} separately reports 134 slow-suite
passes and two-trial smoke runs of all 154 cells. The top-level driver was
parsed, not executed, in the reviewed CI job. The full study remains unrun.

\section*{7. Next step within the agreed scope}
Keep this framework study fixed. First check grid
resolution and retained MCMC draws on selected sparse and heterogeneous
states; document convergence and acceptable factor discrepancies. Then,
\emph{after approval}, run a small timing/operating-characteristic pilot and
choose the production replication count from the required precision.
The earlier three-CPU-hour estimate predates the current 154-cell design;
re-estimate production cost from a pilot of these settings.

Hyperprior sensitivity remains pending. Partial exchangeability, new entry
events, delays, drift, adaptive allocation and the older comparator families
belong to subsequent work. The background paper by Baas, ter Schure and
van Rosmalen (2026), \href{https://arxiv.org/abs/2605.28653}{arXiv:2605.28653},
motivates e-value trial design; its design-optimal comparators are deferred
here. Implementation decisions, review evidence and limitations are recorded
in \texttt{Additional\_files/}.
\end{document}
